\chapter*{Заключение}
\addcontentsline{toc}{chapter}{Заключение}
\label{ch:conclusion}

В рамках настоящей работы решена задача построения визуомоторной политики клонирования поведения для манипулятора, удовлетворяющей требованиям технического задания, и проведено экспериментальное сравнение двух её реализаций.

Базовая политика построена на основе самообучаемого зрительного энкодера DINOv2 и каузального декодерного трансформера со стохастической головой в виде смеси гауссиан. Реализация выполнена поверх библиотеки robomimic и допускает обучение на подмножестве Proficient-Human набора robomimic~v0.1; на обучающем распределении данных достигнута успешность $\text{Success Rate} = 0{,}95$, что существенно превосходит требуемый порог $0{,}75$.

Вторая реализация расширяет базовую архитектуру модулем пространственно-блочного векторного квантования признаков DINOv2 с EMA-обновлением кодовой книги и параллельной ветвью реконструкции исходного кадра по квантованной карте. Совместная оптимизация политики, commitment-потери и попиксельной реконструкции обеспечивает согласованное обучение зрительной части и управляющего трансформера. На обучающем распределении данных алгоритм с квантованием даёт $\text{Success Rate} = 0{,}96$, сопоставимо с базовой политикой.

Проведённое сравнение показало следующее. На обучающем распределении базовый алгоритм и алгоритм с квантованием достигают сопоставимо высоких значений успешности ($0{,}95$ и $0{,}96$ соответственно) и оба с большим запасом превосходят требуемый технический порог $0{,}75$. При доменном сдвиге по цветовой раскладке сцены базовый алгоритм демонстрирует существенное падение качества -- до $0{,}618 \pm 0{,}045$ -- тогда как вариант с VQ-VAE сохраняет уровень успешности $0{,}916 \pm 0{,}013$, фактически неотличимый от исходного. При этом вспомогательные эксперименты показывают, что само по себе квантование, без сопутствующей реконструкционной ветви, не даёт аналогичного прироста: ключевую роль играет именно требование к кодовой книге сохранять структуру изображения, формируемое потерей $\mathcal{L}_{\text{IR}}$. Тем самым экспериментально подтверждается гипотеза, положенная в основу расширения из Раздела~\ref{sec:vq-impl}: пространственно-блочное квантование в сочетании с реконструкцией выступает как механизм формирования доменно-инвариантных визуальных представлений и обеспечивает устойчивость визуомоторной политики к изменениям цветовой палитры сцены без какой-либо дополнительной тренировки или целевой доменной адаптации.

Среди перспективных направлений дальнейшей работы можно выделить систематическое исследование поведения алгоритма при других типах доменных сдвигов (изменение освещения, текстур, геометрии сцены), расширение экспериментов на более сложные задачи набора robomimic помимо Lift, а также изучение влияния размера кодовой книги и веса реконструкционного слагаемого на баланс между качеством политики и устойчивостью представлений.

\endinput
